\section{对偶函数}

在这一章，我们考虑\cref{def:优化问题}给出的标准形式的优化问题
\begin{Equation}[优化问题]
    \begin{array}{lll}
        \minimize   &f_0(\vb{x})\\
        \subto      &f_i(\vb{x})\leq 0, &i=1,\cdots,m\\
                    &h_i(\vb{x})=0,     &i=1,\cdots,p
    \end{array}
\end{Equation}

其中，优化变量$\vb{x}\in\R^n$，优化变量的定义域记为$\mathcal{D}$，最优值记为$p^{*}$。

\subsection{拉格朗日函数}

\begin{BoxDefinition}[拉格朗日函数]
    拉格朗日函数（Lagrangian Function）$\mathcal{L}: \R^n\times\R^m\times\R^p\to\R$，定义为
    \begin{Equation}[拉格朗日函数]
        \mathcal{L}(\vb{x},\vb*{\lambda},\vb*{\nu})=f_0(\vb*{x})+\Sum[i=1][m]\lambda_if_i(\vb{x})=\Sum[i=1][p]\nu_i h_i(\vb{x})
    \end{Equation}
    其中$\lambda_i$和$\nu_i$分别称为$f_i(\vb{x})\leq 0$和$h_i(\vb{x})=0$的拉格朗日乘数（Lagrange Multiplier）。
\end{BoxDefinition}

\subsection{拉格朗日对偶函数}

\begin{BoxDefinition}[拉格朗日对偶函数]
    拉格朗日对偶函数（Lagrange Dual Function）$g: \R^m\times\R^p\to\R$，定义为
    \begin{Equation}[拉格朗日对偶函数]
        g(\vb*{\lambda},\vb*{\nu})=\inf_{\vb{x}\in\mathcal{D}}\mathcal{L}(\vb{x},\vb*{\lambda},\vb*{\nu})
    \end{Equation}
\end{BoxDefinition}

简而言之，拉格朗日对偶$g(\vb*{\lambda},\vb*{\nu})$是拉格朗日$\mathcal{L}(\vb{x},\vb*{\lambda},\vb*{\nu})$关于$\vb{x}\in\mathcal{D}$的下确界。

\begin{BoxProperty}[拉格朗日对偶的下界性]
    拉格朗日对偶在$\vb*{\lambda}\succeq\vb*{0}$时，总是不高于原问题的最优值$p^{*}$
    \begin{Equation}[拉格朗日对偶的下界性]
        \vb*{\lambda}\succeq\vb{0}\qquad g(\vb*{\lambda},\vb*{\nu})\leq p^{*}
    \end{Equation}
\end{BoxProperty}

\begin{Proof}[\cref{ppt:拉格朗日对偶的下界性}]
    根据\cref{def:拉格朗日对偶函数}，设$x^{*}$为最优点，$\mathcal{L}(\vb{x},\vb*{\lambda},\vb*{\nu})$在$\vb{x}\in\mathcal{D}$上的下确界一定小于等于$\mathcal{L}(\vb{x}^{*},\vb*{\lambda},\vb*{\nu})$
    \begin{Equation}[拉格朗日对偶的下界性1]
        g(\vb*{\lambda},\vb*{\nu})=\inf_{\vb{x}\in\mathcal{D}}\mathcal{L}(\vb{x},\vb*{\lambda},\vb*{\nu})\leq\mathcal{L}(\vb{x}^{*},\vb*{\lambda},\vb*{\nu})
    \end{Equation}

    若$\lambda_i\geq 0$，考虑到$f_i(\vb{x}^{*})\leq 0$和$h_i(\vb{x}^{*})=0$
    \begin{Equation}[拉格朗日对偶的下界性2]
        \mathcal{L}(\vb{x}^{*},\vb*{\lambda},\vb*{\nu})=f_0(\vb*{x}^{*})+\Sum[i=1][m]\lambda_if_i(\vb{x}^{*})=\Sum[i=1][p]\nu_i h_i(\vb{x}^{*})\leq f_0(\vb{x}^{*})=p^{*}
    \end{Equation}

    因此就有$g(\vb*{\lambda},\vb*{\nu})\leq p^{*}$成立。
\end{Proof}

